\chapter{Datasets}

Only two datasets are used in the final experiments: RIRD and Wongnai. They serve different
roles because they contain different kinds of supervision.

\section{RIRD: English restaurant retrieval benchmark}
The Reviewed-Item Retrieval Dataset (RIRD) contains 29,168 English reviews from 50 Toronto
restaurants, with between 406 and 2,281 reviews per restaurant. It provides 100 natural-language
preference queries and relevance judgements from five annotators for all restaurant--query
pairs \citep{rird2023}.

RIRD is the main evaluation dataset because it contains the three elements required for
restaurant-level retrieval: review text, restaurant identifiers and query-specific relevance
labels. We remove duplicate judgement rows, aggregate the five annotator fields into binary and
graded relevance, and preserve restaurant and review identifiers through preprocessing. Graded
labels support nDCG, while binary labels support MAP, precision and recall. Since every query
is judged against the same 50 restaurants, models can be compared query by query using paired
tests.

\section{Wongnai: Thai sentence-level case study}
The Wongnai dataset contains approximately 40,000 Thai reviews with text and a 0--4 star
rating. It is sourced from the Hugging Face dataset repository
\citep{wongnai2020}. The available data does not contain restaurant identifiers or
query-specific relevance labels. We therefore do not use Wongnai for restaurant ranking.

We sample 10,000 reviews with seed 42, producing 93,810 Thai sentences. The project creates
24 frozen Thai taste queries covering individual tastes, combinations and negations. BM25 and
dense retrieval produce a blind pool of 452 query--sentence pairs for human assessment. The
star rating is retained as metadata only: a high overall rating cannot prove that a sentence
answers a particular request such as ``I want something sour''. This makes the Thai component
a sentence-level transfer case study, not a second restaurant-ranking benchmark.

\section{Why the other proposed datasets were not used}
The proposal also considered Thailand Foods and the Thai subset of M-ABSA. They are not counted
as experimental datasets in this final report because neither provides the query--restaurant
relevance structure needed for the central evaluation.

Thailand Foods contains dish names and food attributes, so it could support a future food-term
gazetteer or query normalisation layer. It does not contain review evidence, restaurant
identifiers or relevance judgements for our queries. M-ABSA provides aspect--sentiment tuples,
which are useful for Thai sentiment research, but its labels do not measure whether a review
sentence satisfies a restaurant-search query. Using either resource as a ranking ground truth
would require constructing labels or features and would risk evaluating the system against
assumptions introduced by the researchers themselves. We therefore keep both datasets out of
the reported model and evaluation.

\section{Data provenance and limitations}
RIRD and Wongnai are external datasets and are cited in the bibliography. The figures in this
report are generated from our saved experiment outputs, not copied from an external source.
Review text is subjective and may be outdated or contradictory; the system presents it as
retrieval evidence rather than verified fact. RIRD is limited to 50 restaurants, while Wongnai
cannot support restaurant-level evaluation without additional metadata.

\begin{table}[H]
\centering\small
\caption{The two datasets used in the final study.}
\label{tab:datasets}
\begin{tabularx}{\textwidth}{@{}lLLL@{}}
\toprule
Dataset & Content & Role & Main limitation\\
\midrule
RIRD & 29,168 English reviews, 50 restaurants and 100 labelled queries & Restaurant ranking, explanation and baseline evaluation & Toronto restaurants only\\
Wongnai & Approximately 40,000 Thai reviews with 0--4 star ratings & Sentence-level Thai taste transfer & No restaurant IDs or query labels\\
\bottomrule
\end{tabularx}
\end{table}
